\section{Heights and integral frames}
\label{ht:section}

The endomorphisms used in the local equations must satisfy two different
size estimates. Their coordinates may have height proportional to the
height of the target, whereas their positive local logarithms near the
fixed center must have only logarithmic growth in that height. We first
choose small endomorphisms, then construct integral lattices to control
their finite-place coordinates, and finally compare these lattices and
the archimedean Hodge metrics with fixed rational frames.

\subsection{Height conventions and small endomorphisms}

\begin{convention}\label{ht:conventions}
Fix the smooth curve $S$, its smooth projective completion $\overline S$,
the abelian scheme $\A\to S$, and the center $t_0$ supplied by
\Cref{geo:setup,inc:markings}. Enlarge a fixed number field $K$ to contain
all these data and the finitely many fixed charts, levels and center
endomorphisms. Write
\[
 d=\max\{2,[K(t):K]\},\qquad
 h=h_{[t_0]}(t)+c_0\geq 2.
\]
Here $h_{[t_0]}$ is an absolute logarithmic Weil height. The divisor
$[t_0]$ is ample on $\overline S$, and $c_0$ is a fixed constant.
The letters $C$ and $\kappa$ below denote positive constants depending
only on the fixed data; they may be increased in successive estimates.

Choose rational frames of $H_{\rm dR}=H_1^{\rm dR}(\A/S)$ and of
$\cP_{\rm dR}$ that are regular and invertible on an affine neighborhood
of $t_0$. Choose a rational parameter $x$ with a simple zero at $t_0$,
and put $a=x(t)$. Throughout, a sufficiently small disk means the
actual branch about $t_0$, not all points with small $|x|$.
Remove the finitely many poles, frame singularities, and unwanted zeros
of the rational functions being used. Perform the same choices for the
finitely many conjugate data. Such finite omissions are covered by
\Cref{geo:setup}.

For a number field $L$ containing the varying data, put
$n_v=[L_v:\Q_v]/[L:\Q]$ and use the absolute values extending the
usual real or $p$-adic absolute values on $\Q$. In particular,
\[
 h_{\rm aff}(z_1,\ldots,z_r)
   =\sum_v n_v\log\max\{1,|z_1|_v,\ldots,|z_r|_v\}.
\]
At a real or complex place the local degree is understood in the usual
sense. The sum of the archimedean weights is one, and the weights above
any rational prime sum to one. When passing to a finite extension, the
weights of places above a given place sum to its old weight. This also
allows computations after finite local extensions. Weighted sums below
use this convention, including when only a subcollection of embeddings
and places is retained.
\end{convention}

Every fixed rational map from $\overline S$ to a projective space has
height at most $Ch$ wherever it is defined. Indeed its pulled-back
hyperplane divisor is dominated, for height inequalities, by a sufficiently
large multiple of the ample divisor $[t_0]$: their difference can be
chosen ample and its height is bounded below. This proves in particular
$h(a)\leq Ch$ and the same estimate for all fixed base coordinates.

\begin{proposition}\label{ht:endomorphisms}
Let $t$ be a target with
$\End_{\Qbar}(\A_t)\otimes\Q\simeq\Mat_2(\Q)$, and let
$\mathcal O_t=\End_{\Qbar}(\A_t)$. For the principal Rosati involution
$\dagger$, set
\[
 Q_t(f)=\Tr_{H_1}(ff^\dagger).
\]
Then
\begin{equation}\label{ht:small-bounds}
 h_F(\A_t)\leq Ch,\qquad
 |\disc\mathcal O_t|\leq C(dh)^\kappa.
\end{equation}
The discriminant here is that of the full geometric order for the trace
of the regular representation of $\Mat_2(\Q)$.
There are integral endomorphisms $f_1,\ldots,f_4$ spanning
$\mathcal O_t\otimes\Q$ and an independent labeled pair
$\mathbf P=(P_1,P_2)$ spanning its Rosati-symmetric trace-free plane,
all with
\begin{equation}\label{ht:rosati-small}
 Q_t(f_i),\ Q_t(P_j)\leq C(dh)^\kappa.
\end{equation}
All geometric endomorphisms, and hence this pair and its de Rham
coordinates, are defined over an extension $L/K(t)$ of bounded degree.
The bound depends only on the dimension and the fixed data.
\end{proposition}

\begin{proof}
After a fixed finite theta-level cover, a symmetric ample line bundle
and its theta data give an algebraic theta-null map. On that cover use
the pullback of $[t_0]$ as the ample divisor; its height at a lift is
$h_{[t_0]}(t)+O(1)$. The theta-null projective height is therefore at
most $Ch$ by the preceding height comparison. Pazuki's
comparison of theta height and stable Faltings height
\cite[Corollary~1.3(1)]{Pazuki2012} gives $h_F(\A_t)\leq Ch$.
The logarithmic error in that comparison is absorbed into this upper
bound. The cover has fixed degree, so using it does not impose a bound
on the varying degree $d$.

We spell out the field-degree dependence in the endomorphism estimate.
The Proposition of \cite[\S4]{MasserWustholz1994}, with its induction
parameter $m=n$ and with principal polarization, gives integral
endomorphisms spanning the geometric endomorphism algebra with Rosati
quadratic sizes at most
\[
 c\max\{2,[L:\Q],h_F(\A_t)\}^{\kappa_2},
\]
where $c$ and $\kappa_2$ depend only on the dimension $n=2$.
The degree dependence is exposed in that proposition and in
\cite[\S5, equation~(5.3) and Lemma~5.1]{MasserWustholz1994}; it is not
an unspecified constant allowed to depend arbitrarily on $[L:\Q]$.
The estimates apply to arbitrary abelian varieties and thus include the
split algebra $\Mat_2(\Q)$.

Here one can take $[L:K(t)]$ bounded. The continuous Galois action on
the rank-four lattice $\mathcal O_t$ has finite image. To see finiteness,
choose an integral basis; each basis endomorphism is defined over a
finite extension, so an open Galois subgroup fixes the entire basis.
Finite subgroups of $\GL_4(\Z)$ have bounded order. One elementary
bound follows by reduction modulo $3$: its kernel has no nontrivial
torsion. In fact a torsion element in that kernel has a prime-order
power; writing it as $1+3^sA$, with some entry of $A$ a $3$-adic unit,
the binomial expansion rules out order prime to $3$, and also rules out
order $3$ by comparing the valuations of the linear and higher terms.
Thus reduction embeds the finite image in $\GL_4(\mathbf F_3)$.
Taking the fixed field of its kernel proves the assertion about $L$.
Since $K$ is fixed, $[L:\Q]\leq Cd$, which gives the generator bound
in \eqref{ht:rosati-small}.

For completeness, the generator bound controls the full-order
discriminant without requiring the chosen generators to be an integral
basis. Put
\[
 b(f,g)=\Tr_{H_1}(fg),\qquad q(f,g)=b(f,g^\dagger).
\]
The form $q$ is positive definite and integral. Principal polarization
makes $\dagger$ an automorphism of the integral order. Consequently,
in an integral basis, the matrices of $b$ and $q$ differ by right
multiplication by an integral matrix of determinant $\pm1$. Their
determinants have the same absolute value. Select four independent
small generators and let $\Lambda$ be their span over $\Z$.
Hadamard's inequality gives
\[
 \det(q|_\Lambda)\leq\prod_{i=1}^4 Q_t(f_i)
       \leq C(dh)^\kappa,
 \qquad
 \det(q|_{\mathcal O_t})
       =\frac{\det(q|_\Lambda)}{[\mathcal O_t:\Lambda]^2}.
\]
On $\Mat_2(\Q)$ the rank-four homology representation is the sum of
two copies of its standard module. Thus
$\Tr_{H_1}(f)=2\tr_2(f)$, which also equals the trace of the regular
representation. This proves the discriminant convention and its bound.

The $q$-orthogonal projection onto the symmetric trace-free subspace is
\[
 \pi(f)=\frac{f+f^\dagger}{2}
             -\frac{\Tr_{H_1}(f)}4\id.
\]
The two projections commute, and neither increases the $q$-norm.
Moreover
\begin{equation}\label{ht:integral-projection}
 4\pi(f)=2(f+f^\dagger)-\Tr_{H_1}(f)\id\in\mathcal O_t.
\end{equation}
By \Cref{geo:five-space}, the image has dimension two. The images of
the chosen rationally spanning generators span it. Select two
independent nonzero images and use \eqref{ht:integral-projection} to
obtain $P_1,P_2$. Their quadratic sizes increase by at most $16$.
All these operations take place over $L$.
\end{proof}

Fix once and for all an analogous integral pair
$\mathbf B=(B_1,B_2)$ at $t_0$. Neither $\mathbf B$ nor $\mathbf P$
is required to be an integral basis of its plane. The reductions used
in \Cref{inc:markings} concern the saturated plane lattices.

\subsection{A bundle whose trait lattices preserve endomorphisms}

To bound finite-place coordinates, we need a lattice preserved by every
integral endomorphism, including at fibers with bad reduction. We construct
one from the polarization biextension. Its Lie bundle agrees with de Rham
homology on the abelian characteristic-zero open and remains locally free
on the semiabelian model.

\begin{proposition}\label{ht:integral-model}
Let $T$ be a regular noetherian integral scheme, $U\subset T$ a dense
open, and $\mathcal G\to T$ a smooth separated semiabelian scheme of
relative dimension two with geometrically integral fibers. Suppose that
$\mathcal G_U$ is an abelian scheme with principal polarization.
The Poincar\'e bundle, viewed on
$\mathcal G_U\times_U\mathcal G_U$ using that polarization, has a
unique extension as a normalized biextension $\mathscr B$ on
$\mathcal G\times_T\mathcal G$.

This biextension defines a smooth vector-group extension
\begin{equation}\label{ht:vector-extension}
 0\longrightarrow\omega_{\mathcal G/T}
 \longrightarrow\mathcal E\longrightarrow\mathcal G
 \longrightarrow0,
 \qquad \omega_{\mathcal G/T}=e^*\Omega^1_{\mathcal G/T},
\end{equation}
where a locally free module denotes the additive group of its sections.
The vector bundle $\mathscr H=\Lie(\mathcal E)$ has rank four and,
on the characteristic-zero abelian open, identifies functorially with
$H_1^{\rm dR}(\mathcal G_U/U)$.

For any trait $\Spec R\to T$ whose generic point belongs to that open,
and any integral endomorphism $f$ of its generic abelian variety, the
action of $f$ preserves the lattice $\mathscr H_R$. This remains true
after every finite extension of traits on which $f$ is defined.
\end{proposition}

\begin{proof}
We first extend the normalized biextension, then construct its vector
extension, and finally show that endomorphisms act integrally on the
resulting Lie bundle. Every power $\mathcal G^r$ is
regular. A line bundle on $(\mathcal G^r)_U$ has a rational section;
the closure of its divisor is a Weil divisor on $\mathcal G^r$ and is
Cartier by regularity. This extends the line bundle.

Every prime divisor of $\mathcal G^r$ supported over $T\setminus U$
is the reduced pullback of a unique prime divisor of $T$. Indeed
flatness implies that the image of a codimension-one point has
codimension at most one. It cannot be the generic point of $T$.
The fiber above a codimension-one point is geometrically integral,
so its generic point is unique; smoothness makes the pullback divisor
reduced. It follows that two extensions, with their generic
identification, differ by a divisor pulled back from $T$.

Start with any extension $\mathscr B_0$ of the Poincar\'e bundle and
write $p:\mathcal G^2\to T$. Replacing it by
\[
 \mathscr B=\mathscr B_0\otimes
         p^*((e,e)^*\mathscr B_0)^{-1}
\]
gives the prescribed trivialization at the simultaneous origin.
The generic axis rigidifications extend uniquely: the divisor of each
rational rigidification is pulled back from $T$, and restriction to
the origin, where it is already the identity, makes that divisor zero.
For the same reason, the two generic biextension isomorphisms extend
on $\mathcal G^3$. For example, the ratio line bundle in
\[
 (m,1)^*\mathscr B\simeq
          p_{13}^*\mathscr B\otimes p_{23}^*\mathscr B
\]
has a rational trivialization whose vertical divisor is detected at
$(e,e,e)$. It is therefore zero. A rational section of a line bundle
on a regular scheme with zero divisor is an everywhere invertible
section. Associativity, compatibility of the two laws, and the unit
identities hold on the dense open, hence everywhere. The same argument
proves uniqueness of a normalized extension. This is the divisorial
argument used over a trait in
\cite[Proposition~6.7.1]{EdixhovenLido2023}; the argument just given
establishes the stated version over $T$.

Let $N$ be the first infinitesimal neighborhood of the identity in
the second copy of $\mathcal G$. Its ideal is square-zero, and
$\cO_N=\cO_T\oplus\omega_{\mathcal G/T}$ as an augmented
$\cO_T$-module. Over a test scheme $V\to T$, an element of
$\mathcal E(V)$ is a point $g\in\mathcal G(V)$ together with a
trivialization of $\mathscr B|_{\{g\}\times N_V}$ extending its
given value at the origin. Such trivializations exist locally on $V$.
Two of them differ by an element of
$1+\omega_{\mathcal G/T}\otimes\cO_V$, whose multiplication is
addition because the ideal is square-zero. The resulting torsor is
represented explicitly as follows. If
$q:\mathcal G\times_TN\to\mathcal G$ is the finite flat projection
and $\pi:\mathcal G\to T$ is the structure map, axis normalization
gives an exact sequence of vector bundles
\[
 0\longrightarrow\pi^*\omega_{\mathcal G/T}
 \longrightarrow q_*(\mathscr B|_{\mathcal G\times_TN})
 \longrightarrow\cO_{\mathcal G}\longrightarrow0.
\]
The inverse image of the section $1$ is the required affine torsor.
Its points are actual sections of this fixed bundle; no quotient by
automorphisms of semiabelian extensions is taken. The first
biextension law supplies its group law. Its fiber over $e$ is the
vector group $\omega_{\mathcal G/T}$, with its canonical origin.
This proves \eqref{ht:vector-extension}, including smoothness and
compatibility with arbitrary base change. Taking Lie algebras gives
an exact sequence of vector bundles of ranks $2,4,2$.

Over the abelian open, this is the universal vector extension of
$\mathcal G_U$, expressed as infinitesimally rigidified line bundles
on its dual; see \cite[I, \S2.6 and \S4]{MazurMessing1974} and
\cite[Proposition~2.2]{Cais2010}. In cohomological notation the
classical Lie realization is
\[
 \Lie E(A)=H^1_{\rm dR}(A^\vee).
\]
The Poincar\'e de Rham pairing identifies the right side with
$H_1^{\rm dR}(A)$, with the usual Tate line understood over the base
field. This is the covariant realization: a map $f:A\to A$ acts via
$(f^\vee)^*$ on $H^1_{\rm dR}(A^\vee)$ and via $f_*$ on homology.
The Hodge filtration and functoriality in this statement are made
explicit in \cite[Proposition~5.1 and Corollary~5.2]{Cais2010}.
Only this characteristic-zero identification is needed. On a toric
special fiber we retain the constructed rank-four Lie bundle; no
universal property or integral de Rham duality is asserted there.

Now pull back to a trait. The normal-base homomorphism extension
theorem \cite[Chapter~I, Proposition~2.7]{FaltingsChai1990},
\cite[Proposition~3.3.1.5]{Lan2013} applies directly to the
semiabelian scheme $\mathcal G_R$. It extends both $f$ and
$f^\dagger=\lambda^{-1}f^\vee\lambda$ to $\mathcal G_R$;
principal polarization is what makes $f^\dagger$ integral.
The generic normalized biextension identity
\begin{equation}\label{ht:adjoint-biextension}
 (f,1)^*\mathscr B\simeq(1,f^\dagger)^*\mathscr B
\end{equation}
extends over $R$ by the same divisor argument. Pulling an
infinitesimal rigidification back along $f^\dagger|_N$ therefore
sends the fiber of $\mathcal E$ over $g$ to its fiber over $f(g)$.
It defines an $R$-homomorphism $\mathcal E\to\mathcal E$.
On the generic fiber it is the covariant map just described.
Its differential preserves $\Lie(\mathcal E_R)$, proving the
integrality assertion. After a finite trait extension the base-changed
group scheme is again semiabelian and the new trait is normal.
Apply the same homomorphism extension theorem and the same normalized
biextension argument there. Compatibility of the constructed torsor
with base change then proves the final assertion.
\end{proof}

We apply \Cref{ht:integral-model} to each of the two regular integral
toroidal models used in \Cref{inc:torus}, after forgetting the other
auxiliary level. For $i=1,2$ put
$R_i=\cO_K[1/\ell_i]$; the roots of unity required by the levels
already belong to $K$. Over $R_i$, take a proper projective graph
model $\mathscr X_i$ of $\overline S$ mapping to the level-$\ell_i$
toroidal compactification, and pull the rank-four bundle back to
$\mathscr X_i$. The graph model need not be regular: local freeness
survives pullback, and the proposition is applied on the regular
toroidal base. A point $t$ extends to $\mathscr X_i$ at every finite
place whose residue characteristic is different from $\ell_i$.
Since $\ell_1\ne\ell_2$, these two constructions supply a lattice
at every finite place.

For the height calculation, extend each graph model separately as
follows. Choose a projective embedding
$\mathscr X_i\hookrightarrow\mathbf P^{M_i}_{R_i}$ and let
$\overline{\mathscr X}_i$ be its reduced scheme-theoretic closure in
$\mathbf P^{M_i}_{\cO_K}$. This is an integral proper projective
model with
$\overline{\mathscr X}_i\times_{\cO_K}R_i=\mathscr X_i$.
The bundle is used only on this $R_i$-open. Fixed line-bundle
denominators on $\overline{\mathscr X}_i$, constructed below, provide
global height bounds even though the bundle is unavailable over
$\ell_i$.

\subsection{Global coordinates and the sharper estimate near the center}

\begin{lemma}\label{ht:frame-costs}
For each finite place let $\mathscr H_t$ be the lattice obtained from
an available model above. There are nonnegative local costs
$\delta_v(t)$, depending only on the fixed frames and models, such
that the coefficients in the chosen $\cP_{\rm dR}$ frame of any
integral endomorphism $f\in\cP_{{\rm dR},t}$ satisfy
\begin{equation}\label{ht:finite-frame-bound}
 \log^+\|f\|_{v,\mathrm{coord}}\leq\delta_v(t),
 \qquad
 \sum_{v\nmid\infty}n_v\delta_v(t)\leq Ch.
\end{equation}
At sufficiently close places about a potentially-good center, the
costs have bounded weighted sum. Outside a fixed finite set of rational
primes they can be taken to be zero on those disks. These assertions
are uniform over finite extensions and over the finitely many
conjugate centers.
\end{lemma}

\begin{proof}
Choose a rational frame $e_1,\ldots,e_4$ for $H_{\rm dR}$ and its
dual frame, and include the rational coefficient maps from
$\End(H_{\rm dR})$ to the chosen coordinates on $\cP_{\rm dR}$.
For each $i$, these form a finite list of rational sections of vector
bundles on $\mathscr X_i$. Extend each of those bundles to a coherent
sheaf $\mathscr F$ on $\overline{\mathscr X}_i$. Such an extension
exists by clearing powers of $\ell_i$ in finite sets of local
generators and relations on a finite affine cover. Removing its
torsion preserves its restriction to $\mathscr X_i$.
No local freeness of $\mathscr F$ over $\ell_i$ is needed.

Let $r$ be one of the rational sections. The local functions $b$ for
which $br$ belongs to $\mathscr F$ form a coherent nonzero ideal
sheaf $\mathscr I_r$. Its stalk at the characteristic-zero center
$t_0$ is the unit ideal, because the chosen frames and coefficient maps
are regular there. For a sufficiently high power of the ample line bundle
$\mathscr L_i=\cO_{\mathbf P^{M_i}}(1)|_{\overline{\mathscr X}_i}$,
$\mathscr I_r\otimes\mathscr L_i^{\otimes m}$ is generated by global
sections. We may therefore choose a section nonzero at $t_0$. Viewed
as a section of $\mathscr L_i^{\otimes m}$, it gives a fixed
$s\in H^0(\overline{\mathscr X}_i,\mathscr L_i^{\otimes m})$
such that $sr$ is a regular section of
$\mathscr F\otimes\mathscr L_i^{\otimes m}$.
The zero divisor of $s$ is an effective Cartier divisor, since
$\overline{\mathscr X}_i$ is integral. Make these choices for the
finite list of sections on both models and omit their finitely many
generic-fiber zeros on the curve; the center remains in the resulting
open. The same choices can be made for each of the fixed conjugate data.

At places of residue characteristic different from $\ell_i$, equip
the vector bundles with the lattice norms from $\mathscr X_i$.
Use the model metric of $\mathscr L_i$ at \emph{all} finite places,
including those above $\ell_i$, and continuous metrics at infinity.
For a point off the zero divisor put
$\lambda_{s,v}(t)=-\log\|s(t)\|_v$.
Regularity of $sr$ bounds the positive logarithmic norm of $r(t)$
by $\lambda_{s,v}(t)$ at the available places. Every finite-place
$\lambda_{s,v}(t)$ is nonnegative, because $s$ is an integral
section on the proper model $\overline{\mathscr X}_i$.
Moreover its archimedean terms are uniformly bounded below, since
$\|s\|$ is bounded above on the compact complex curve. Consequently
\[
 \sum_{\substack{v\nmid\infty\\v\nmid\ell_i}}
       n_v\lambda_{s,v}(t)
 \leq\sum_{v\nmid\infty}n_v\lambda_{s,v}(t)
 \leq h_{\mathscr L_i^{\otimes m}}(t)+C
 \leq Ch.
\]
The height in this display uses the generic-fiber line bundle and
the fixed metrics just specified. In particular no vector-bundle
lattice over $\ell_i$ has been used to bound the partial sum.
Adding the finitely many denominator bounds gives the required
global frame cost.

If $f$ preserves $\mathscr H_t$, then its operator norm for that
lattice is at most one. Its matrix coefficients are evaluations
$e_i^*(fe_j)$ and are bounded by the product of the norms of $e_i^*$
and $e_j$. The fixed coordinate maps on $\cP$ contribute their own
denominator costs. Adding these finitely many terms proves
\eqref{ht:finite-frame-bound}. Applying the argument for both models
and adding their bounds covers all finite places.

For the refinement, outside finitely many rational primes the center
extends to the abelian interior and the rational frames, their duals,
and the coordinate maps are integral and invertible throughout its
formal neighborhood. Their costs are zero there. At one of the
remaining potentially-good center places, first make a fixed finite
local extension giving good reduction. The center then lies in the
interior of a prime-to-characteristic integral model. On a sufficiently
small disk in its actual branch, the family remains in that interior
with the same specialization. Since the rational frames are regular
and invertible at the characteristic-zero center, their matrices and
inverse matrices are bounded on a smaller closed disk. This follows
by evaluating their finitely many regular coordinate expressions;
the disk radius and the bounds depend only on the center and the
expressions. It remains valid for points over arbitrarily ramified
extensions. Each such rational prime contributes a fixed constant,
because its total place weight is at most one. There are only
finitely many such primes and conjugate centers.
\end{proof}

At an archimedean place use the Hodge metric on $H_1(\A_t,\C)$
defined by its principal polarization. For an actual endomorphism,
Rosati adjunction is adjunction for this metric, and therefore
\begin{equation}\label{ht:hodge-operator}
 \|f\|_{\mathrm{op,Hdg}}\leq Q_t(f)^{1/2}.
\end{equation}
The same trace integer computes $Q_t(f)$ after every field embedding,
so this estimate and \eqref{ht:rosati-small} are simultaneous at all
archimedean places.

Here is the needed control of the frame norms. At a puncture with
parameter $u$, make a fixed finite cover $u=z^b$ to make monodromy
unipotent, using \cite[Monodromy Theorem~(6.1)]{Schmid1973}.
Regularity of the Gauss--Manin connection follows from
\cite[II, Proposition~6.14 and Theorem~7.9]{Deligne1970}, applied to
the smooth proper family and the trivial coefficient connection.
Choose the \emph{canonical unipotent extension}, with nilpotent
residue, from \cite[II, \S5, Proposition~5.2]{Deligne1970}.
If $T$ is the unipotent monodromy, $N=\log T$, and $\mathbf v$
is a flat frame on the universal cover, this extension is locally
framed by the single-valued sections obtained by cancelling monodromy:
\[
 \mathbf w(z)=
     \exp\!\left(-\frac{\log z}{2\pi\mathrm i}N\right)
     \mathbf v(z).
\]
Both this change-of-frame matrix and its inverse are polynomials in
$\log z$, since $N$ is nilpotent. A different holomorphic frame of
this same extension differs from $\mathbf w$ by a holomorphic
invertible matrix whose norm and inverse norm are bounded on a smaller
disk. The algebraic rational frame differs from such a frame by
meromorphic matrices with fixed finite pole orders; the agreement
with the algebraic de Rham structure is also stated in
\cite[Theorem~(4.13)(a),(b)]{Schmid1973}.
Schmid's estimates for a
polarized variation, applied also to the dual variation, bound the
Hodge norms of flat frames and their duals by powers of
$1+|\log|z||$; see
\cite[Theorem~(6.6$'$), p.~252]{Schmid1973}.
It follows that, for a fixed exponent $M$, all the frame norms and
dual norms needed here are at most
\begin{equation}\label{ht:arch-growth}
 C|u|^{-M}(1+|\log|u||)^M
\end{equation}
near each missing point. Finitely many sectors suffice, and the bound
is independent of the target. The meromorphic coordinate maps from
endomorphisms to $\cP$ have the same type of bound. After taking
positive logarithms, \eqref{ht:arch-growth} is at most a fixed
multiple of a local height of the puncture plus a constant. Summing
at the archimedean places is at most $Ch$, by the same effective
divisor comparison as above. On a fixed small disk about an interior
center these frame norms are simply bounded.

\begin{proposition}\label{ht:coordinates}
The labeled pair $\mathbf P$ of \Cref{ht:endomorphisms}, expressed
in the fixed rational $\cP_{\rm dR}$ frame, satisfies
\begin{equation}\label{ht:global-coordinates}
 h_{\rm aff}(\mathbf P)\leq C(h+\log d).
\end{equation}
Let $V$ be any subcollection of sufficiently close places in the
actual center branches, with bad-center places excluded as in
\Cref{inc:bad-disks}. Then
\begin{equation}\label{ht:near-coordinates}
 \sum_{v\in V}n_v\log\max\{1,\|P_1\|_{v,\mathrm{coord}},
                                  \|P_2\|_{v,\mathrm{coord}}\}
       \leq C(1+\log h+\log d).
\end{equation}
Both estimates hold for simultaneously conjugated data. They also
hold for a fixed number of endomorphisms $F_j$ with
$NF_j\in\mathcal O_t$ for a fixed positive integer $N$ and
$Q_t(F_j)\leq C(dh)^\kappa$.
If $\phi:\A_t\to\A_{t'}$ is a polarized isogeny of one of the
fixed multipliers in \Cref{geo:hecke-list}, write $I_t=\Ad(\phi)$,
so $I_t(D)=\phi D\phi^{-1}$. In particular the estimates hold for
the transported pair $\widehat{\mathbf P}=I_t^{-1}\mathbf P'$.
The added choices in this last assertion have bounded relative field
degree.
\end{proposition}

\begin{proof}
At finite places use \Cref{ht:frame-costs}. At infinity combine
\eqref{ht:hodge-operator} with the frame and dual-frame bounds.
The logarithm of the Rosati size is at most
$C(1+\log d+\log h)$. The global frame cost is at most $Ch$, and
$\log h\leq h$; this proves \eqref{ht:global-coordinates}.
Near the good center branches, the finite-place cost is bounded and
the archimedean frames have bounded norms. This gives
\eqref{ht:near-coordinates}. The construction is invariant under
simultaneous conjugation, and all constants may be maximized over
the fixed finite collection of conjugate models and frames.

For $NF_j$ integral, the finite-place calculation acquires at most
$\log^+|N|_v^{-1}$. Its total weighted sum is $\log N$, and its
sum on any subcollection is no larger. The Hodge calculation applies
directly to $F_j$ with its asserted Rosati bound. This proves the
bounded-denominator assertion.

Finally let $\phi:(\A_t,\lambda_t)\to
(\A_{t'},\lambda_{t'})$ have fixed multiplier $m$, so
$\phi^\vee\lambda_{t'}\phi=m\lambda_t$ and
$\phi^{-1}=m^{-1}\phi^\dagger$. Its conjugation map on
endomorphisms is a Rosati isometry. For $P_j'$ integral, the
transported endomorphism is
\[
 \widehat P_j=\phi^{-1}P_j'\phi
             =m^{-1}\phi^\dagger P_j'\phi.
\]
Thus $m\widehat P_j$ is integral and
$Q_t(\widehat P_j)=Q_{t'}(P_j')$. Apply the result just proved;
this avoids estimating the coordinate height of a matrix of $\phi$.
There are uniformly boundedly many such $\phi$ for fixed $m$:
its kernel is contained in $\A_t[m]$, and, for each kernel, the
polarized identifications of the quotient with $\A_{t'}$ form a
torsor under a finite polarized automorphism group. The latter has
bounded order, since its faithful action on integral homology is a
finite subgroup of $\GL_4(\Z)$. The kernel choices and hence the
isogeny choices are bounded in number by a constant depending only
on $m$. Their Galois orbits have that bound, so choosing $\phi$
requires only a bounded field extension of the compositum of the
two fields of data. There are only finitely many allowed $m$.
\end{proof}

All estimates in this section allow arbitrary $d$ and $h$.
In particular, \eqref{ht:near-coordinates} was obtained from integral
actions and local frame bounds before establishing any upper bound
for $h$ in terms of $d$.
